% Generated draft: requires mathematical and source review.
\documentclass[11pt]{article}
\usepackage[margin=25mm]{geometry}
\usepackage{fontspec}
\setmainfont{Latin Modern Roman}
\usepackage{amsmath,amssymb,mathtools,unicode-math}
\setmathfont{Latin Modern Math}
\usepackage{enumitem,needspace,xurl,hyperref}
\setlength{\parindent}{0pt}
\setlength{\parskip}{5pt}
\setlength{\emergencystretch}{3em}
\newcommand{\N}{\mathbb{N}}
\newcommand{\Z}{\mathbb{Z}}
\newcommand{\Q}{\mathbb{Q}}
\newcommand{\R}{\mathbb{R}}
\newcommand{\C}{\mathbb{C}}
\newcommand{\F}{\mathbb{F}}
\newcommand{\A}{\mathbb{A}}
\newcommand{\PP}{\mathbb{P}}
\newcommand{\E}{\mathbb{E}}
\newcommand{\eps}{\varepsilon}
\newcommand{\dd}{\,\mathrm{d}}
\begin{document}
\section*{UnsolvedMath ID: 801039}
\textbf{Source problem code:} catalogue entry
\subsection{Definitions and mathematical statement}
Let \(X\) be a standard Borel space and let \(E\subseteq X\times X\) be a Borel equivalence relation with countable equivalence classes. This uses the conventional standard Borel setting for a countable Borel equivalence relation. Write \([x]_E=\{y\in X:xEy\}\). The relation \(E\) is \emph{hyperfinite} if there are Borel equivalence relations \(F_n\) on \(X\), for \(n\in\N\), such that every \(F_n\)-class is finite, \(F_n\subseteq F_{n+1}\), and
\[
E=\bigcup_{n\in\N}F_n.
\]
The source calls \(E\) \emph{hyperfinite-over-hyperfinite}, abbreviated \emph{hf/hf}, if there is a Borel partial order \(\leq\) on \(X\) satisfying
\[
\forall x,y\in X\quad (x\leq y\Longrightarrow xEy)
\]
and, for every \(x\in X\), the restriction of \(\leq\) to \([x]_E\) has order type a suborder of \(\Z^2\).

\textbf{Unresolved qualification.} The supplied record does not specify an order on \(\Z^2\); its opening phrase, ``Analogous to the above,'' refers to context not included in the record. Consequently the hf/hf condition cannot be made unambiguous from this record alone. If \(\preceq_*\) denotes the missing intended order, then under the usual order-embedding interpretation of suborder the classwise requirement is
\[
\forall x\in X\ \exists f_x:[x]_E\longrightarrow\Z^2\quad
\left[ f_x\text{ is injective and }\forall u,v\in[x]_E\;
\bigl(u\leq v\Longleftrightarrow f_x(u)\preceq_* f_x(v)\bigr)\right].
\]
Here \(\preceq_*\) is a placeholder for missing source information, not a chosen order. This classwise existence condition does not impose a Borel or uniform choice of the maps \(f_x\).

\textbf{Question (the hf/hf problem).} With the intended order and meaning of suborder restored, is every hf/hf countable Borel equivalence relation hyperfinite? Equivalently, for every such \(X\) and \(E\), does the existence of the stated Borel partial order imply the existence of an increasing sequence \((F_n)_{n\in\N}\) as above?
\subsection{Short English statement}
The hf/hf problem asks whether a countable Borel equivalence relation is hyperfinite whenever its classes admit a compatible Borel partial order whose classwise order types occur inside the intended ordered \(\Z^2\). The supplied record omits the order on \(\Z^2\), so the precise hypothesis requires clarification.
\subsection{Sources}
\begin{itemize}
\item Supplied local UnsolvedMath registry record: permanent ID \(801039\), exact source code \texttt{catalogue entry}, title \emph{Open Problems in Mathematical Logic --- Problem 51 (The hf/hf problem)}. The record lists Barmpalias, Gao et al. as proposers, with proposed year \(2026\), and labels the problem open; this is supplied registry status, not an independently verified assessment.
\item Barmpalias, Gao et al., \emph{Open Problems in Mathematical Logic} (2026), Problem 51, \url{https://arxiv.org/abs/2608.26628}. The supplied provenance reports access on 2026-10-16 and extraction by \texttt{arxiv\_latex}; these are supplied provenance details, not a claim of independent source access. The supplied attribution/context is ``The hf/hf problem.''
\end{itemize}
\end{document}
